% GENERATED FILE. Edit canonical lessons, then run npm run generate:books.
% canonical-source-hash: fnv1a64:624e493acc47133c
% canonical-lessons: IT-C372-supporre, IT-C372-ipotesi, IT-C372-presumere, IT-C372-plausibile, IT-C372-presupposto

\chapter{Supponiamo}
\label{ch:it-supponiamo}
\hlchaptermodality{372}

\begin{chapteropening}
\textbf{By the end of this chapter:} \emph{I can make a hypothesis: suppose or presume something, say it is plausible and name what it rests on.}

Complete the last lesson of chapter 372: I can make a hypothesis: suppose or presume something, say it is plausible and name what it rests on.
\end{chapteropening}

\section[supporre]{supporre — to suppose, to assume}

\label{lesson:IT-C372-supporre}



\begin{quote}
Before the new one: say the Italian for fundamental, then the Italian for the premise.
\end{quote}

\subsection*{You'll want to know: supporre}
\textbf{supporre} — \textquotedblleft{}to suppose, to assume\textquotedblright{}.

\begin{quote}
Supponiamo che domani piova: che cosa facciamo?
\end{quote}

\begin{quote}
\emph{Let us suppose it rains tomorrow: what do we do?}
\end{quote}

Conjugated like \emph{porre}: \emph{suppongo, supponi}; past participle \emph{supposto}. \emph{Supponiamo che} takes the subjunctive.

\begin{cousinweb}
\emph{supporre} comes from Latin \emph{suppōnere} \textquotedblleft{}to put under, to substitute\textquotedblright{} (Lewis \& Short), built from \emph{sub} \textquotedblleft{}under\textquotedblright{} and \emph{pōnere} \textquotedblleft{}to put\textquotedblright{}: an idea placed underneath the rest.
\end{cousinweb}

\begin{grammarlens}[title={What you've built}]
One more word for telling, explaining and describing at B1.
\end{grammarlens}

\subsection*{Your turn}
\begin{itemize}
\raggedright
  \item \emph{Say it:} \emph{supporre}
  \item \emph{Say it:} \emph{supporre}, once more
  \item \emph{Say it:} a sentence of your own with \emph{supporre}
  \item \emph{Recall:} say the Italian for fundamental, then the Italian for the premise, then say \emph{supporre} again
\end{itemize}

\begin{tcolorbox}[breakable,colback=teal!4,colframe=teal!35!black,title={Before you move on}]
What does supporre mean? (\textquotedblleft{}To suppose, to assume\textquotedblright{}.) Say it once more.
\end{tcolorbox}
\section[l'ipotesi]{l'ipotesi — the hypothesis, the possibility}

\label{lesson:IT-C372-ipotesi}



\begin{quote}
Before the new one: say the Italian for the premise, then the Italian for to suppose.
\end{quote}

\subsection*{You'll want to know: l'ipotesi}
\textbf{l'ipotesi} — \textquotedblleft{}the hypothesis, the possibility\textquotedblright{}.

\begin{quote}
Nell'ipotesi che il treno sia in ritardo, prendiamo un taxi.
\end{quote}

\begin{quote}
\emph{If the train happens to be late, we will take a taxi.}
\end{quote}

The same in the plural: \emph{le ipotesi}. \emph{Nell'ipotesi che} takes the subjunctive.

\begin{cousinweb}
\emph{ipotesi} is a learned borrowing from Late Latin \emph{hypothesis}, from Greek \emph{hypóthesis} \textquotedblleft{}a placing under, a supposition\textquotedblright{}, built from \emph{hypó} \textquotedblleft{}under\textquotedblright{} and \emph{thésis} \textquotedblleft{}a placing\textquotedblright{}, the word behind Latin \emph{thesis} \textquotedblleft{}a proposition\textquotedblright{} (Lewis \& Short).
\end{cousinweb}

\begin{grammarlens}[title={What you've built}]
One more word for telling, explaining and describing at B1.
\end{grammarlens}

\subsection*{Your turn}
\begin{itemize}
\raggedright
  \item \emph{Say it:} \emph{l'ipotesi}
  \item \emph{Say it:} \emph{l'ipotesi}, once more
  \item \emph{Say it:} a sentence of your own with \emph{l'ipotesi}
  \item \emph{Recall:} say the Italian for the premise, then the Italian for to suppose, then say \emph{l'ipotesi} again
\end{itemize}

\begin{tcolorbox}[breakable,colback=teal!4,colframe=teal!35!black,title={Before you move on}]
What does l'ipotesi mean? (\textquotedblleft{}The hypothesis, the possibility\textquotedblright{}.) Say it once more.
\end{tcolorbox}
\section[presumere]{presumere — to presume, to take for granted}

\label{lesson:IT-C372-presumere}



\begin{quote}
Before the new one: say the Italian for to suppose, then the Italian for the hypothesis.
\end{quote}

\subsection*{You'll want to know: presumere}
\textbf{presumere} — \textquotedblleft{}to presume, to take for granted\textquotedblright{}.

\begin{quote}
Presumo che tu sia stanco dopo il viaggio.
\end{quote}

\begin{quote}
\emph{I presume you are tired after the journey.}
\end{quote}

Past participle \emph{presunto}. \emph{Presuntuoso} means arrogant.

\begin{cousinweb}
\emph{presumere} is a learned borrowing from Latin \emph{praesūmere} \textquotedblleft{}to take beforehand, to take for granted\textquotedblright{} (Lewis \& Short), built from \emph{prae} \textquotedblleft{}before\textquotedblright{} and \emph{sūmere} \textquotedblleft{}to take\textquotedblright{}.
\end{cousinweb}

\begin{grammarlens}[title={What you've built}]
One more word for telling, explaining and describing at B1.
\end{grammarlens}

\subsection*{Your turn}
\begin{itemize}
\raggedright
  \item \emph{Say it:} \emph{presumere}
  \item \emph{Say it:} \emph{presumere}, once more
  \item \emph{Say it:} a sentence of your own with \emph{presumere}
  \item \emph{Recall:} say the Italian for to suppose, then the Italian for the hypothesis, then say \emph{presumere} again
\end{itemize}

\begin{tcolorbox}[breakable,colback=teal!4,colframe=teal!35!black,title={Before you move on}]
What does presumere mean? (\textquotedblleft{}To presume, to take for granted\textquotedblright{}.) Say it once more.
\end{tcolorbox}
\section[plausibile]{plausibile — plausible, believable}

\label{lesson:IT-C372-plausibile}



\begin{quote}
Before the new one: say the Italian for the hypothesis, then the Italian for to presume.
\end{quote}

\subsection*{You'll want to know: plausibile}
\textbf{plausibile} — \textquotedblleft{}plausible, believable\textquotedblright{}.

\begin{quote}
La tua spiegazione è plausibile, ma non mi convince.
\end{quote}

\begin{quote}
\emph{Your explanation is plausible, but it does not convince me.}
\end{quote}

One form for both genders.

\begin{cousinweb}
\emph{plausibile} is a learned borrowing from Latin \emph{plausibilis} \textquotedblleft{}deserving applause, acceptable\textquotedblright{} (Lewis \& Short), from \emph{plaudere} \textquotedblleft{}to clap the hands\textquotedblright{}.
\end{cousinweb}

\begin{grammarlens}[title={What you've built}]
One more word for telling, explaining and describing at B1.
\end{grammarlens}

\subsection*{Your turn}
\begin{itemize}
\raggedright
  \item \emph{Say it:} \emph{plausibile}
  \item \emph{Say it:} \emph{plausibile}, once more
  \item \emph{Say it:} a sentence of your own with \emph{plausibile}
  \item \emph{Recall:} say the Italian for the hypothesis, then the Italian for to presume, then say \emph{plausibile} again
\end{itemize}

\begin{tcolorbox}[breakable,colback=teal!4,colframe=teal!35!black,title={Before you move on}]
What does plausibile mean? (\textquotedblleft{}Plausible, believable\textquotedblright{}.) Say it once more.
\end{tcolorbox}
\section[il presupposto]{il presupposto — the precondition, the assumption}

\label{lesson:IT-C372-presupposto}



\begin{quote}
Before the new one: say the Italian for to presume, then the Italian for plausible.
\end{quote}

\subsection*{You'll want to know: il presupposto}
\textbf{il presupposto} — \textquotedblleft{}the precondition, the assumption\textquotedblright{}.

\begin{quote}
Il presupposto è che tutti siano d'accordo.
\end{quote}

\begin{quote}
\emph{The assumption is that everyone agrees.}
\end{quote}

\emph{Partire dal presupposto che} — to start from the assumption that.

\begin{cousinweb}
\emph{presupposto} is the past participle of \emph{presupporre} \textquotedblleft{}to presuppose\textquotedblright{}, built from \emph{pre-} \textquotedblleft{}before\textquotedblright{} and \emph{supporre}, from Latin \emph{suppōnere} \textquotedblleft{}to put under\textquotedblright{} (Lewis \& Short), on \emph{pōnere} \textquotedblleft{}to put\textquotedblright{}.
\end{cousinweb}

\begin{grammarlens}[title={What you've built}]
One more word for telling, explaining and describing at B1.
\end{grammarlens}

\subsection*{Your turn}
\begin{itemize}
\raggedright
  \item \emph{Say it:} \emph{il presupposto}
  \item \emph{Say it:} \emph{il presupposto}, once more
  \item \emph{Say it:} a sentence of your own with \emph{il presupposto}
  \item \emph{Recall:} say the Italian for to presume, then the Italian for plausible, then say \emph{il presupposto} again
\end{itemize}

\begin{tcolorbox}[breakable,colback=teal!4,colframe=teal!35!black,title={Before you move on}]
What does il presupposto mean? (\textquotedblleft{}The precondition, the assumption\textquotedblright{}.) Say it once more.
\end{tcolorbox}
